%mscS --> seminar
%msc --> Thesis

% برای گزارش سمینار mscS  و برای پایان‌نامه ,msc را انتخاب کنید
\documentclass[oneside,openany,mscS]{SBU-Thesis}

% برای وارد کردن بسته‌هایی که نیاز است قبل از bidi وارد شوند، آنها را در فایل confix.tex قرار دهید. سایر بسته‌های مورد نیاز را در اینجا می‌توانید وارد کنید:

\begin{document}
	%عنوان پایان‌نامه
	\title{به‌کارگیری رویکرد‌های تطابق معنایی در بازیابی اشیا}
	
	%رشته
	\subject{مهندسی کامپیوتر}
	
	%گرایش
	\field{نرم‌افزار}
	
	%نام
	\name{آرش}
	
	%نام خانوادگی
	\surname{درگاهی نوبری}
	
	%استاد راهنما
	\firstsupervisor{دکتر محمود نشاطی}
	
	%استاد راهنمای دوم (در صورت وجود، در غیر این صورت این خط را حذف کنید)
%	\secondsupervisor{دکتر محمود نشاطی}

	%استاد مشاور (در صورت وجود، در غیر این صورت این خط را حذف کنید)
%	\firstadvisor{دکتر محمود نشاطی}

%تاریخ انجام پایان‌نامه
	\thesisdate{زمستان ۱۳۹۷}
	
	%نام دانشکده
	\faculty{دانشکده مهندسی و علوم کامپیوتر}
	
	%نام دانشگاه
	\university{دانشگاه شهید بهشتی}
	
	%اسامی برای صفحات داوران: در هر مورد، در صورت عدم وجود، کل خط را حذف کنید
	\davaranSupervisor{نام و نام‌خانوادگی}
	\davaranSecondSupervisor{نام و نام‌خانوادگی}
	\davaranAdvisor{نام و نام‌خانوادگی}
	\davaranInternal{نام و نام‌خانوادگی}
	\davaranExternal{نام و نام‌خانوادگی}
	\davaranAssignee{نام و نام‌خانوادگی}
%	\davaranDate{۱۳۹۷/۱۱/۷} % چنانچه این خط را حذف کنید با فضای خالی جایگزین می‌شود

	%مقادیر انگلیسی برای صفحه آخر
	\latinuniversity{Shahid Beheshti University}
	\latinfaculty{Faculty of Computer Science \& Engineering}
	\latintitle{Semantic Matching Methods on Object Retrieval}
	\latinname{Arash}
	\latinsurname{Dargahi Nobari}
	\latinthesisdate{2019}
	\firstlatinsupervisor{Dr. Mahmood Neshati}
	%	\secondlatinsupervisor{Dr. Mahmood Neshati}
	%	\firstlatinadvisor{Dr. Mahmood Neshati}
	% چکیده به فارسی
	\fa-abstract{
		لورم ایپسوم ( به انگلیسی \lr{lorem ipsum} ) متنی بی مفهوم است که تشکیل شده از کلمات معنی دار یا بی معنی کنار هم. کاربر با دیدن متن لورم ایپسوم تصور میکند متنی که در صفحه مشاهده میکند این متن واقعی و مربوط به توضیحات صفحه مورد نظر است واقعی است. حالا سوال اینجاست که این متن « لورم ایپسوم » به چه دردی میخورد و اساسا برای چه منظور و هدفی ساخته شده است؟ پیش از بوجود آمدن لورم ایپسوم ، طراحان وب سایت در پروژه های وب سایت و طراحان کرافیک در پروژه های طراحی کاتولوگ ، بروشور ، پوستر و ... همواره با این مشکل مواجه بودند که صفحات پروژه خود را پیش از آنکه متن اصلی توسط کارفرما ارائه گردد و در صفحه مورد نظر قرار گیرد چگونه پر کنند؟؟ اکثر طراحان با نوشتن یک جمله مانند «این یک متن نمونه است» ویا «توضیحات در این بخش قرار خواهند گرفت» و کپی آن به تعداد زیاد یک یا چند پاراگراف متن میساختند که تمامی متن ها و کلمات ، جملات و پاراگراف ها تکراری بود و از این رو منظره خوبی برای بیننده نداشت و ضمنا به هیچ وجه واقعی به نظر نمیرسید تا بتواند شکل و شمایل تمام شده پروژه را نشان دهد. از این رو متنی ساخته شد که با دو کلمه ( به فارسی : لورم ایپسوم ) آغاز میشد وبا همین نام در بین طراحان وب و گرافیک شناخته و به سرعت محبوب شد. وب سایت های سازنده لورم ایپسوم میتوانند هر تعداد کلمه و پاراگراف که بخواهید به صوورت تکراری یا غیر تکراری برایتان بسازند و تحویلتان بدهند تا از آنها در پروژه هایتان استفاده کنید. ( لورم ایپسوم فارسی) متن های لورم ایپسوم را به زبان فارسی و علاوه بر زبان فارسی به انگلیسی ، عربی ، ترکی استانبولی و ... برایتان میسازد. زبان های دیگر نیز رفته رفته به بانک اطلاعاتی لورم ایسپوم فارسی اضافه خواهند شد.  
	}
	%کلمات کلیدی:
	\keywords{واژه۱، واژه۲، واژه۳}
	%%%%%%%%%%%%%%%%%%%%%%%%%%%%%%%%%
	
	
\firstPage %ساخت صفحه اول پایان‌نامه 
\davaranPage % ساخت صفحه امضا داوران

%%%%%%ساخت صفحه تقدیر و تشکر%%%%%%%%%%
{
	\newpage
	\thispagestyle{plain}
	\noindent
	\large{\textbf{با سپاس و قدردانی از}}
	
	\noindent
	پدران و مادرانی که خود را فدای تربیت فرزاندان خود کردند و\\
	اساتید و معلمانی که در تمام دوران زندگی، راهنمای جانسوز ما بودند.
	
	
	\vspace{14cm}	
	آوردن این صفحه اختیاریست.
	
	\pagebreak
}
%%%%%%%%%%%%%%%%%%%%%%%%%%%%%%%%

\rightsPage % ساخت صفحه حقوق پایان‌نامه
\copyRightPage % ساخت صفحه تعهدنامه

%%%%%%ساخت صفحه تقدیم به %%%%%%%%%%
{
	\newpage
	\thispagestyle{plain}
	\large{\textbf{تقدیم به}}
	
	\begin{center}
		رهجويان علم و فناوری و دوست‌داران علم و دانش
	\end{center}
	
	\vspace{14cm}	
	آوردن این صفحه اختیاریست.
	
	\pagebreak
}


%%%%%%%%%%%%%%%%%%%%%
\abstractPage % ساخت صفحه چکیده


\tableofcontents % فهرست مطالب
\listoffigures \newpage % فهرست تصاویر
\listoftables \newpage % فهرست جداول

		

%%%%%%%%%%%%%%%%%%%%%%%%%%%

\cchapter{مقدمه}
در این بخش یک معرفی اجمالی از پایان نامه آورده می‌شود. تمام متنی که در این قالب آمده صرفاً جهت نمایش قالب و نشان دادن قابلیت‌های آن است.
این قالب با استفاده از بسته 
\lr{XePersian}
ساخته شده و از تمام قابلیت‌های رایج \lr{ LaTeX} می‌توان در نگارش آن استفاده کرد. و برای استفاده تمام دانشجویان دانشگاه شهید بهشتی منتشر شده است.
در بخش‌‌های بعدی به نحوه استفاده از این قالب می‌پردازیم. در ادامه متنی بی معنا جهت پر شدن حجم صفحه آورده می‌شود:

لورم ایپسوم ( به انگلیسی \lr{lorem ipsum} ) متنی بی مفهوم است که تشکیل شده از کلمات معنی دار یا بی معنی کنار هم. کاربر با دیدن متن لورم ایپسوم تصور میکند متنی که در صفحه مشاهده میکند این متن واقعی و مربوط به توضیحات صفحه مورد نظر است واقعی است. حالا سوال اینجاست که این متن « لورم ایپسوم » به چه دردی میخورد و اساسا برای چه منظور و هدفی ساخته شده است؟ پیش از بوجود آمدن لورم ایپسوم ، طراحان وب سایت در پروژه های وب سایت و طراحان کرافیک در پروژه های طراحی کاتولوگ ، بروشور ، پوستر و ... همواره با این مشکل مواجه بودند که صفحات پروژه خود را پیش از آنکه متن اصلی توسط کارفرما ارائه گردد و در صفحه مورد نظر قرار گیرد چگونه پر کنند؟؟ اکثر طراحان با نوشتن یک جمله مانند «این یک متن نمونه است» ویا «توضیحات در این بخش قرار خواهند گرفت» و کپی آن به تعداد زیاد یک یا چند پاراگراف متن میساختند که تمامی متن ها و کلمات ، جملات و پاراگراف ها تکراری بود و از این رو منظره خوبی برای بیننده نداشت و ضمنا به هیچ وجه واقعی به نظر نمیرسید تا بتواند شکل و شمایل تمام شده پروژه را نشان دهد. از این رو متنی ساخته شد که با دو کلمه ( به فارسی : لورم ایپسوم ) آغاز میشد وبا همین نام در بین طراحان وب و گرافیک شناخته و به سرعت محبوب شد. وب سایت های سازنده لورم ایپسوم میتوانند هر تعداد کلمه و پاراگراف که بخواهید به صوورت تکراری یا غیر تکراری برایتان بسازند و تحویلتان بدهند تا از آنها در پروژه هایتان استفاده کنید. ( لورم ایپسوم فارسی) متن های لورم ایپسوم را به زبان فارسی و علاوه بر زبان فارسی به انگلیسی ، عربی ، ترکی استانبولی و ... برایتان میسازد. زبان های دیگر نیز رفته رفته به بانک اطلاعاتی لورم ایسپوم فارسی اضافه خواهند شد.  لورم ایپسوم ( به انگلیسی \lr{lorem ipsum} ) متنی بی مفهوم است که تشکیل شده از کلمات معنی دار یا بی معنی کنار هم. کاربر با دیدن متن لورم ایپسوم تصور میکند متنی که در صفحه مشاهده میکند این متن واقعی و مربوط به توضیحات صفحه مورد نظر است واقعی است. حالا سوال اینجاست که این متن « لورم ایپسوم » به چه دردی میخورد و اساسا برای چه منظور و هدفی ساخته شده است؟ پیش از بوجود آمدن لورم ایپسوم ، طراحان وب سایت در پروژه های وب سایت و طراحان کرافیک در پروژه های طراحی کاتولوگ ، بروشور ، پوستر و ... همواره با این مشکل مواجه بودند که صفحات پروژه خود را پیش از آنکه متن اصلی توسط کارفرما ارائه گردد و در صفحه مورد نظر قرار گیرد چگونه پر کنند؟؟ اکثر طراحان با نوشتن یک جمله مانند «این یک متن نمونه است» ویا «توضیحات در این بخش قرار خواهند گرفت» و کپی آن به تعداد زیاد یک یا چند پاراگراف متن میساختند که تمامی متن ها و کلمات ، جملات و پاراگراف ها تکراری بود و از این رو منظره خوبی برای بیننده نداشت و ضمنا به هیچ وجه واقعی به نظر نمیرسید تا بتواند شکل و شمایل تمام شده پروژه را نشان دهد. از این رو متنی ساخته شد که با دو کلمه ( به فارسی : لورم ایپسوم ) آغاز میشد وبا همین نام در بین طراحان وب و گرافیک شناخته و به سرعت محبوب شد. وب سایت های سازنده لورم ایپسوم میتوانند هر تعداد کلمه و پاراگراف که بخواهید به صوورت تکراری یا غیر تکراری برایتان بسازند و تحویلتان بدهند تا از آنها در پروژه هایتان استفاده کنید. ( لورم ایپسوم فارسی) متن های لورم ایپسوم را به زبان فارسی و علاوه بر زبان فارسی به انگلیسی ، عربی ، ترکی استانبولی و ... برایتان میسازد. زبان های دیگر نیز رفته رفته به بانک اطلاعاتی لورم ایسپوم فارسی اضافه خواهند شد.  

\cchapter{معرفی قابلیت‌های قالب}
در این بخش با آوردن یک متن ساده به نمایش ظاهر و ساختار قالب و همچنین معرفی برخی دستورات لازم برای کار با آن قالب پرداخته می‌شود. این توضیحات بسیار مختصر بوده و صرفاً برای معرفی قالب می‌باشد و چنانچه با \lr{LaTeX} آشنایی ندارید، بهتر است پیش از آغاز تدوین پایان‌نامه مختصری در مورد نحوه کار با \lr{LaTeX} مطالعه بفرمایید.
\section{نگارش}
رعایت تمامی اصول نگارش در هنگام تدوین پایان‌نامه الزامیست، بسیاری از نکات نگارشی توسط قالب رعایت می‌شوند. در ادامه این بخش به معرفی برخی دستورات کاربردی برای این کار پرداخته می‌شود.
\\
به طور پیش‌‌فرض هر پاراگراف به صورت خودکار با فاصله  از کنار آغاز می‌شود.  چنانچه در حالت خاصی، نیاز به حذف این فاصله باشد می‌توانید از دستور 
\LTR{\lr{\textbackslash noindent}} \RTL{}\noindent
استفاده کنید.\\
برای نوشتن متون انگلیسی لازم است آن‌ها را داخل تگ
 \lr{\textbackslash lr\{\}}
 قرار دهید. به عنوان مثال با نوشتن 
 \lr{\textbackslash lr\{word\}}
 کلمه \lr{word} به درستی داخل متن قرار می‌گیرد.
 \par
 با نوشتن \textbackslash\textbackslash می‌توانید خط جدیدی را آغاز کنید و نوشتن \textbackslash par نیز باعث ایجاد یک پاراگراف جدید خواهد شد. با قرار دادن متن در داخل تگ \lr{\textbackslash textbf\{\}} متن به صورت \textbf{ضخیم} و با قرار دادن نوشته در داخل تگ \lr{\textbackslash textit\{\}} نوشته \textit{کج} خواهد شد. امکان استفاده هم‌زمان از این تگ‌ها نیز وجود دارد. برای آشنایی بیشتر با دستورات، پیشنهاد می‌شود به آموزش‌های \lr{LaTeX} مراجعه کنید.
 \par

 پانوشت‌‌ها یکی از بخش‌های اصلی در هر نوشته‌ای می‌باشد. در این قالب شماره‌های پانوشت در هر صفحه، مجدداً از ۱ آغاز می‌شود.  با نوشتن
\LTR{\lr{\textbackslash LTRfootnote\{footnote\}}} \RTL{} \noindent
  می‌توان یک پانوشت لاتین\LTRfootnote{footnote} اضافه نمود.
 \par
برای افزودن پانوشت‌های فارسی\RTLfootnote{یک پانوشت فارسی} نیز از دستور زیر استفاده می‌شود.
\LTR{\lr{\textbackslash RTLfootnote\{پانوشت\}}} \RTL{} \noindent
نمونه‌های افزودن پانوشت نیز در همین قسمت وجود دارد.






\section{بخش‌بندی}
برای ساخت یک فصل جدید کافیست از دستور
\LTR{\lr{\textbackslash cchapter\{عنوان\}}} \RTL{} \noindent
استفاده شود. با نوشتن این دستور به صورت خودکار یک فصل جدید اضافه شده و عنوان آن در یک صفحه مجزا قرار می‌گیرد.
هر بخش می‌تواند شامل تعدادی \lr{Section} باشد. شماره‌های آن مانند آنچه در بالا نیز مشاهده می‌کنید با . از یک‌دیگر جدا شده و به صورت خودکار شماره‌گزاری  شده و به فهرست اضافه می‌شوند. کافیست برای ساخت بخش دستور 
\LTR{\lr{\textbackslash section\{عنوان\}}} \RTL{}\noindent
را وارد کنید. همچنین با دستورات زیر می‌توانید زیربخش و حتی زیر زیر بخش، ایجاد کنید.
\LTR{\lr{\textbackslash subsection\{عنوان\}}} \RTL{}
\LTR{\lr{\textbackslash subsubsection\{عنوان\}}} \RTL{}\noindent
به عنوان مثال بخش زیر را در نظر بگیرید :(متون این زیربخش‌ها بی‌معنا و برای پر کردن صفحه می‌باشد.)

\subsection{یک زیر بخش}
اگر زیر بخش‌ها به سطح سوم برسند شماره‌گذاری نمی‌شوند ولی در فهرست مطالب قرار می‌گیرند، به عنوان مثال:
\subsubsection{زیر بخشی در زیر بخش}
متن
\subsubsection{ دومین زیربخش در زیربخش}
متنی دیگر

\subsection{زیر بخشی دیگر}
لورم ایپسوم ( به انگلیسی \lr{lorem ipsum} ) متنی بی مفهوم است که تشکیل شده از کلمات معنی دار یا بی معنی کنار هم. کاربر با دیدن متن لورم ایپسوم تصور میکند متنی که در صفحه مشاهده میکند این متن واقعی و مربوط به توضیحات صفحه مورد نظر است واقعی است. حالا سوال اینجاست که این متن « لورم ایپسوم » به چه دردی میخورد و اساسا برای چه منظور و هدفی ساخته شده است؟ پیش از بوجود آمدن لورم ایپسوم ، طراحان وب سایت در پروژه های وب سایت و طراحان کرافیک در پروژه های طراحی کاتولوگ ، بروشور ، پوستر و ... همواره با این مشکل مواجه بودند که صفحات پروژه خود را پیش از آنکه متن اصلی توسط کارفرما ارائه گردد و در صفحه مورد نظر قرار گیرد چگونه پر کنند؟؟ اکثر طراحان با نوشتن یک جمله مانند «این یک متن نمونه است» ویا «توضیحات در این بخش قرار خواهند گرفت» و کپی آن به تعداد زیاد یک یا چند پاراگراف متن میساختند که تمامی متن ها و کلمات ، جملات و پاراگراف ها تکراری بود و از این رو منظره خوبی برای بیننده نداشت و ضمنا به هیچ وجه واقعی به نظر نمیرسید تا بتواند شکل و شمایل تمام شده پروژه را نشان دهد. از این رو متنی ساخته شد که با دو کلمه ( به فارسی : لورم ایپسوم ) آغاز میشد وبا همین نام در بین طراحان وب و گرافیک شناخته و به سرعت محبوب شد. وب سایت های سازنده لورم ایپسوم میتوانند هر تعداد کلمه و پاراگراف که بخواهید به صوورت تکراری یا غیر تکراری برایتان بسازند و تحویلتان بدهند تا از آنها در پروژه هایتان استفاده کنید. ( لورم ایپسوم فارسی) متن های لورم ایپسوم را به زبان فارسی و علاوه بر زبان فارسی به انگلیسی ، عربی ، ترکی استانبولی و ... برایتان میسازد. زبان های دیگر نیز رفته رفته به بانک اطلاعاتی لورم ایسپوم فارسی اضافه خواهند شد.  لورم ایپسوم ( به انگلیسی \lr{lorem ipsum} ) متنی بی مفهوم است که تشکیل شده از کلمات معنی دار یا بی معنی کنار هم. کاربر با دیدن متن لورم ایپسوم تصور میکند متنی که در صفحه مشاهده میکند این متن واقعی و مربوط به توضیحات صفحه مورد نظر است واقعی است. حالا سوال اینجاست که این متن « لورم ایپسوم » به چه دردی میخورد و اساسا برای چه منظور و هدفی ساخته شده است؟ پیش از بوجود آمدن لورم ایپسوم ، طراحان وب سایت در پروژه های وب سایت و طراحان کرافیک در پروژه های طراحی کاتولوگ ، بروشور ، پوستر و ... همواره با این مشکل مواجه بودند که صفحات پروژه خود را پیش از آنکه متن اصلی توسط کارفرما ارائه گردد و در صفحه مورد نظر قرار گیرد چگونه پر کنند؟؟ اکثر طراحان با نوشتن یک جمله مانند «این یک متن نمونه است» ویا «توضیحات در این بخش قرار خواهند گرفت» و کپی آن به تعداد زیاد یک یا چند پاراگراف متن میساختند که تمامی متن ها و کلمات ، جملات و پاراگراف ها تکراری بود و از این رو منظره خوبی برای بیننده نداشت و ضمنا به هیچ وجه واقعی به نظر نمیرسید تا بتواند شکل و شمایل تمام شده پروژه را نشان دهد. از این رو متنی ساخته شد که با دو کلمه ( به فارسی : لورم ایپسوم ) آغاز میشد وبا همین نام در بین طراحان وب و گرافیک شناخته و به سرعت محبوب شد. وب سایت های سازنده لورم ایپسوم میتوانند هر تعداد کلمه و پاراگراف که بخواهید به صوورت تکراری یا غیر تکراری برایتان بسازند و تحویلتان بدهند تا از آنها در پروژه هایتان استفاده کنید. ( لورم ایپسوم فارسی) متن های لورم ایپسوم را به زبان فارسی و علاوه بر زبان فارسی به انگلیسی ، عربی ، ترکی استانبولی و ... برایتان میسازد. زبان های دیگر نیز رفته رفته به بانک اطلاعاتی لورم ایسپوم فارسی اضافه خواهند شد.  

\section{ارجاعات}
دادن ارجاعات با افزودن \lr{BibTeX} مرجع به فایل \lr{thesis.bib} و سپس آوردن نام آن انجام می‌شود. مانند همین ارجاعی که در این صفحه وجود دارد. ارجاعات می‌تواند به منابع انگلیسی\cite{expret} یا فارسی\cite{Vahedi87} باشد.

\section{نمایه‌ها}
در این بخش به نحوه افزودن فایل‌های خارجی و جداول به متن پرداخته می‌شود.
\subsection{فرمول‌نویسی}
نگارش فرمول‌ها در \lr{LaTeX} به صورتی که مشاهده‌ می‌کنید انجام می‌شود، در ادامه چند فرمول به عنوان نمونه نوشته می‌شود. تمامی قواعد و قوانین فرمول‌نویسی در \lr{LaTeX} بدون نیاز به هیچ‌گونه تغییری در این قالب قابل استفاده می‌باشد.
\begin{equation}
MAP(Q) = \frac{1}{|Q|} \sum_{j = 1}^{|Q|}
{
	\frac{1}{m_j} \sum_{k = 1}^{m_j} {Precision(R_jk)}
}
\label{eq:map}
\end{equation}
\begin{equation}
MI(sa,w) = \sum_{ A_{sa} = 0,1}{
	\sum_{A_{w} = 0,1}{
		p(A_{sa} , A_{w}) \log{\frac{p(A_{sa} , A_{w})}{p(A_{sa}) p(A_{w})}}
	}%inner sum
}%First sum
\label{eq:mi}
\end{equation}
برای ارجاع به فرمول‌ها می‌توان از دستور \lr{\textbackslash ref\{label\}} استفاده کرد. به عنوان مثال: معادله \ref{eq:map} از معادلات مهم است.
\subsection{تصاویر}
برای افزودن تصاویر می‌توانید از فایل‌‌های برداری مانند \lr{PDF} و یا تصاویر پیکسلی مانند \lr{PNG} و \lr{JPG} استفاده کنید. افزودن این تصاویر مانند تصویر زیر انجام می‌شود و به راحتی می‌توان برای آنها توضیح نوشت و به طور خودکار در فهرست تصاویر قرار می‌گیرند. این تصاویر به طور خودکار در مکان مناسب قرار می‌گیرند.
\par
آوردن نام تصاویر در متن نیز مانند فرمول‌ها می‌باشد به عنوان مثال: تصویر \ref{fig:logo} لوگو دانشگاه را نشان می‌دهد.
\begin{figure}[tb]
	\centering
	\includegraphics[width=.6\linewidth]{logo}
	\caption {یک نمونه تصویر}
	\label{fig:logo}
\end{figure}
\par
متنی که در ادامه این بخش آمده فقط برای پرکردن فضا می‌باشد: لورم ایپسوم ( به انگلیسی \lr{lorem ipsum} ) متنی بی مفهوم است که تشکیل شده از کلمات معنی دار یا بی معنی کنار هم. کاربر با دیدن متن لورم ایپسوم تصور میکند متنی که در صفحه مشاهده میکند این متن واقعی و مربوط به توضیحات صفحه مورد نظر است واقعی است. حالا سوال اینجاست که این متن « لورم ایپسوم » به چه دردی میخورد و اساسا برای چه منظور و هدفی ساخته شده است؟ پیش از بوجود آمدن لورم ایپسوم ، طراحان وب سایت در پروژه های وب سایت و طراحان کرافیک در پروژه های طراحی کاتولوگ ، بروشور ، پوستر و ... همواره با این مشکل مواجه بودند که صفحات پروژه خود را پیش از آنکه متن اصلی توسط کارفرما ارائه گردد و در صفحه مورد نظر قرار گیرد چگونه پر کنند؟؟ اکثر طراحان با نوشتن یک جمله مانند «این یک متن نمونه است» ویا «توضیحات در این بخش قرار خواهند گرفت» و کپی آن به تعداد زیاد یک یا چند پاراگراف متن میساختند که تمامی متن ها و کلمات ، جملات و پاراگراف ها تکراری بود و از این رو منظره خوبی برای بیننده نداشت و ضمنا به هیچ وجه واقعی به نظر نمیرسید تا بتواند شکل و شمایل تمام شده پروژه را نشان دهد. از این رو متنی ساخته شد که با دو کلمه ( به فارسی : لورم ایپسوم ) آغاز میشد وبا همین نام در بین طراحان وب و گرافیک شناخته و به سرعت محبوب شد. وب سایت های سازنده لورم ایپسوم میتوانند هر تعداد کلمه و پاراگراف که بخواهید به صوورت تکراری یا غیر تکراری برایتان بسازند و تحویلتان بدهند تا از آنها در پروژه هایتان استفاده کنید. ( لورم ایپسوم فارسی) متن های لورم ایپسوم را به زبان فارسی و علاوه بر زبان فارسی به انگلیسی ، عربی ، ترکی استانبولی و ... برایتان میسازد. زبان های دیگر نیز رفته رفته به بانک اطلاعاتی لورم ایسپوم فارسی اضافه خواهند شد.  لورم ایپسوم ( به انگلیسی \lr{lorem ipsum} ) متنی بی مفهوم است که تشکیل شده از کلمات معنی دار یا بی معنی کنار هم. کاربر با دیدن متن لورم ایپسوم تصور میکند متنی که در صفحه مشاهده میکند این متن واقعی و مربوط به توضیحات صفحه مورد نظر است واقعی است. حالا سوال اینجاست که این متن « لورم ایپسوم » به چه دردی میخورد و اساسا برای چه منظور و هدفی ساخته شده است؟ پیش از بوجود آمدن لورم ایپسوم ، طراحان وب سایت در پروژه های وب سایت و طراحان کرافیک در پروژه های طراحی کاتولوگ ، بروشور ، پوستر و ... همواره با این مشکل مواجه بودند که صفحات پروژه خود را پیش از آنکه متن اصلی توسط کارفرما ارائه گردد و در صفحه مورد نظر قرار گیرد چگونه پر کنند؟؟ اکثر طراحان با نوشتن یک جمله مانند «این یک متن نمونه است» ویا «توضیحات در این بخش قرار خواهند گرفت» و کپی آن به تعداد زیاد یک یا چند پاراگراف متن میساختند که تمامی متن ها و کلمات ، جملات و پاراگراف ها تکراری بود و از این رو منظره خوبی برای بیننده نداشت و ضمنا به هیچ وجه واقعی به نظر نمیرسید تا بتواند شکل و شمایل تمام شده پروژه را نشان دهد. از این رو متنی ساخته شد که با دو کلمه ( به فارسی : لورم ایپسوم ) آغاز میشد وبا همین نام در بین طراحان وب و گرافیک شناخته و به سرعت محبوب شد. وب سایت های سازنده لورم ایپسوم میتوانند هر تعداد کلمه و پاراگراف که بخواهید به صوورت تکراری یا غیر تکراری برایتان بسازند و تحویلتان بدهند تا از آنها در پروژه هایتان استفاده کنید. ( لورم ایپسوم فارسی) متن های لورم ایپسوم را به زبان فارسی و علاوه بر زبان فارسی به انگلیسی ، عربی ، ترکی استانبولی و ... برایتان میسازد. زبان های دیگر نیز رفته رفته به بانک اطلاعاتی لورم ایسپوم فارسی اضافه خواهند شد.  


\subsection{جداول}
ساخت جداول در این قالب مانند بقیه متون نوشته شده توسط \lr{LaTeX} می‌باشد. ابزارهای مختلفی نیز برای ساخت خودکار این جداول موجود می‌باشد که می‌توانید از آنها  برای ساخت جداول پیچیده‌تراستفاده کنید. در ادامه نمونه‌ای جدول آورده می‌شود. برای ارجاع به نام جداول در متن نیز مانند تصاویر عمل می‌شود. به عنوان مثال در ادامه جدول \ref{tab:tch} را مشاهده می‌کنید.

\begin{table}[h]
	\begin{center}
%		\def\arraystretch{2}
		\caption{نمونه‌ای جدول}
		\begin{tabular}{|c|c|c|c|c|}
			\hline
			روش & 
			دقت & صحت & حساسیت & معیار اِف \\
			\hline
			
			روش اول & 
			\% 98/70 & 0/8821 & 0/8817 & 0/8819 \\
			\hline
			
			روش دوم & 
			\% 98/15 & 0/7464 & 0/4772 & 0/6118 \\
			\hline
		\end{tabular}
		\label{tab:tch}
	\end{center}
\end{table}



%%%%%%%%%%%%%%%%%%%%%%%%%%%
	

% مراجع
\newpage
%\onehalfspacing
\bibliographystyle{ieeetr-fa}%{chicago-fa}%{plainnat-fa}%
\bibliography{thesis}


%%%%%%%%%% نمونه‌ی واژه‌نلمه تک ستونه ٪٪٪٪٪٪٪٪٪٪
\baselineskip=.6cm

\chapter*{واژه‌نامه  انگلیسی به  فارسی}\markboth{واژه‌نامه  انگلیسی به  فارسی}{واژه‌نامه  انگلیسی به  فارسی}
%\thispagestyle{fancy}
\addcontentsline{toc}{chapter}{واژه‌نامه  انگلیسی به  فارسی}
\noindent
\englishgloss{Word}{کلمه}
\englishgloss{Word 2}{کلمه ۲}
\englishgloss{Word 3}{کلمه ۳}


\chapter*{واژه‌نامه فارسی به انگلیسی}\markboth{واژه‌نامه فارسی به انگلیسی}{واژه‌نامه فارسی به انگلیسی}
\addcontentsline{toc}{chapter}{واژه‌نامه فارسی به انگلیسی}
%\thispagestyle{fancy}
\noindent
\persiangloss{کلمه ۱}{word 1}
\persiangloss{کلمه ۲}{word 2}
\persiangloss{کلمه ۳}{word 3}

\baselineskip=1cm
%%%%%%%%%%%%%%%%%%%%%%%%%


% چکیده به انگلیسی
\en-abstract
{
	This is Abstract in English.
}

% کلمات کلیدی انگلیسی 
\latinkeywords
{
	Word1, Word2, Word3
}

\latinAbstractPage % ساخت صفحه چکیده به انگلیسی
\latinFirstPage % ساخت صفحه آخر

	
\end{document}